\section{Microphone Arrays: Acoustic Sources Separation}

\subsection{Introduction}
Let us consider the presence of $N$ acoustic sources in an environment, emitting different signal:

\[
s_{1}(n), \ s_{2}(n), \ \cdots, \ s_{N}(n)
\]

Such sources could be, for instance, multiple speakers in a room, musical instruments in a recording session, or environmental noise sources such as machinery and traffic.

We acquire the signals using $M > 1$ microphones, placed at different positions in space.  
Each microphone captures a mixture of the source signals, affected by the acoustic propagation paths (i.e. direct and reflected) and potentially corrupted by noise.
Formally, the microphone signals are denoted as:

\[
y_{1}(n), \ y_{2}(n), \ \cdots, \ y_{M}(n)
\]

Each of them contains a linear combination of the $N$ source signals, shaped by the room impulse responses between each source and microphone.
The \textbf{goal} is to obtain an estimation of the \textbf{original source signals} from the observed microphone signals.  
This process, known as \textbf{acoustic source separation}, is essential for applications such as speech enhancement, meeting transcription, music signal processing and noise suppression in hearing aids.

\subsection{Taxonomy of separation techniques}
Different \textbf{taxonomies} can be adopted to categorise the separation techniques used to tackle the problem of acoustic source separation.  
These taxonomies help identify the \textbf{most suitable approach} depending on the characteristics of the environment, the recording setup and the available prior knowledge.

One relevant aspect is \textbf{how much information} we have about the \textbf{mixing process}, the way in which the signals from the sources combine at the microphones.  
This knowledge can range from having no prior information, as in blind source separation, to knowing the exact characteristics of the mixing system, such as the impulse responses.

Another important factor is whether the environment is \textbf{reverberating or not}.  
In anechoic conditions, signals propagate without reflections, making separation easier.
In reverberant conditions, the microphones capture multiple delayed and attenuated copies of each source, significantly increasing the complexity of the problem.

Finally, the relationship between the \textbf{number of sources} and the \textbf{number of microphones} plays a crucial role.  
If the number of sources is smaller than or equal to the number of microphones, the separation problem is more tractable.  
If there are more sources than microphones, the system is underdetermined: we require more sophisticated processing methods.

\subsubsection{Blind vs non-blind techniques}
The first question to consider in the taxonomy of source separation techniques is: \textbf{how much information} do we have about the \textbf{mixing process}?
The answer to this question determines whether we are working in a \textbf{blind} or \textbf{non-blind} setting.

In \textbf{Blind Source Separation (BSS)}, the only available information consists of the microphone signals and, possibly, some a priori statistical knowledge of the source signals.
The goal is to infer, without direct knowledge of the mixing system, the parameters of a demixing process that can recover the original sources from the observed mixtures.
This is inherently an ill-posed problem and some indeterminacy in the demixing model remains, even after estimation.

On the other hand, \textbf{Non-Blind Source Separation} assumes that we have some knowledge about the channel, typically by means of impulse responses.
With this information, the objective is to invert a known (or partially known) system, which may be convolutional in nature.
This additional knowledge generally simplifies the problem compared to the blind case, although challenges remain, especially in reverberant or noisy environments.

\subsubsection{Instantaneous and convolutive mixing}
The second question in our taxonomy concerns the environment: is it \textbf{reverberant or not}?  
The answer determines whether the mixing is considered \textbf{instantaneous} or \textbf{convolutive}.

In the case of \textbf{instantaneous mixing}, there is no reverberation (or it is negligible).  
The microphone signals can be modelled as \textbf{linear combinations} of the source signals, possibly with complex-valued weights $a_{ij}$:

\[
\begin{bmatrix}
y_{1}(t) \\
\vdots \\
y_{M}(t)
\end{bmatrix}
=
\underbrace{
\begin{bmatrix}
a_{11} & \cdots & a_{1N} \\
\vdots & \ddots & \vdots \\
a_{M1} & \cdots & a_{MN}
\end{bmatrix}
}_{\text{mixing matrix}}
\begin{bmatrix}
s_{1}(t) \\
\vdots \\
s_{N}(t)
\end{bmatrix}
\]

In this scenario, the separation problem reduces to finding and inverting the mixing matrix.

In contrast, \textbf{convolutive mixing} arises in reverberant environments: each microphone signal is a filtered version of the source signals, where the \textbf{filters} $g_{ik}(n)$ correspond to the \textbf{acoustic paths} between the sources and the microphones.  
The system is described by the following expression:

\[
y_{k}(n) =
\sum_{i=1}^{N} \sum_{l=0}^{L-1} g_{ik}(l) \, s_{i}(n - l), 
\qquad k = 1, \cdots, M
\]

The separation task now involves inverting a convolutive system, which is significantly more challenging.

\subsubsection{Underdetermined and overdetermined source separation}
The third question addresses the relationship between the \textbf{number of microphones} and the \textbf{number of sources}: is the number of microphones greater than the number of sources?

When the number of microphones satisfies $ M > N $, the problem is said to be \textbf{overdetermined}.
There are more measurements than unknown signals, which generally makes the separation task easier, as there is enough information to solve for the sources directly.

On the other hand, when $ M < N $, the problem is \textbf{underdetermined}.
There are fewer measurements than sources, making the separation not directly solvable without additional assumptions: it becomes necessary to estimate the mixing model using some a priori knowledge of the \textbf{statistical or structural properties} of the source signals.

\subsection{Example: source separation using binary masks}
Let us introduce a specific scenario: we consider an \textit{underdetermined} problem, with $M = 2$ microphones and $N = 3$ sources.
The task is approached as a \textit{Blind Source Separation problem}, assuming \textit{low reverberation} in the environment.
The sources are \textit{fixed in space} and consist of speech signals.

The idea is to perform the separation in the \textbf{time-frequency domain} using the Short-Time Fourier Transform (STFT).  
The method relies on \textbf{binary masking} of the spectrogram to isolate the contributions of the different sources.

The motivation for this is that signals tend to be \textit{more sparse} in the time-frequency domain than in the time domain.  
This sparsity means that, at many time-frequency points, only one source is dominant.  
Such a property makes it easier to isolate individual sources by applying binary masks to the spectrogram, effectively zeroing out time-frequency regions where a given source is absent or weaker.

\begin{figure}[H]
    \centering
    \includegraphics[width=0.7\linewidth]{Binary masking.png}
    \caption{Binary masking for acoustic source separation: sources (left) and observations (right) spectograms}
\end{figure}

We denote by $Y_1(n,\omega_k)$ and $Y_2(n,\omega_k)$ the microphone signals STFTs and by $S_1(n,\omega_k)$, $S_2(n,\omega_k)$, $S_3(n,\omega_k)$ the source signals STFTs.
An \textbf{assumption} we can make, which is reasonable for speech signals, is that at each \textbf{time-frequency pair} $(n,\omega_k)$, the microphone signals are associated to \textbf{only one source}.
The goal is to find three \textbf{binary masks} $M_1, M_2, M_3$ which we can use to extract the source STFTs from one microphone STFT:

\[
\hat{S}_i(n,\omega_k) = Y_1(n,\omega_k) \cdot M_i(n,\omega_k), \qquad i = 1,2,3
\]

In order to find the solution, we can exploit \textbf{spatial information} to separate the sources.
In absence of reverberation, the microphone signals contain attenuated and delayed versions of the source signals.
For bins associated to the $i$-th source:

\[
\begin{cases}
Y_1(n,\omega_k) = \alpha_{1i} \cdot S_i(n,\omega_k) \ e^{-j\omega_k\tau_{1i}} \\[5pt]
Y_2(n,\omega_k) = \alpha_{2i} \cdot S_i(n,\omega_k) \ e^{-j\omega_k\tau_{2i}}
\end{cases} \qquad i = 1,2,3
\]

Each time-frequency bin associated with the $i$-th source exhibits \textbf{constant time-delay} $T_i$ and \textbf{constant attenuation ratio} $A_i$:

\[
T_i = \tau_{2i} - \tau_{1i},
\qquad
A_i = \frac{\alpha_{2i}}{\alpha_{1i}}
\]

Regarding the methodology, for each time-frequency bin, we first compute the two quantities previously defined.
The \textbf{amplitude ratio} $A(n,\omega_k)$ and the \textbf{time delay} $T(n,\omega_k)$ are:

\[
A(n,\omega_k) = \frac{\abs{Y_2(n,\omega_k)}}{\abs{Y_1(n,\omega_k)}},
\qquad
T(n,\omega_k) = \frac{\angle Y_2(n,\omega_k) - \angle Y_1(n,\omega_k)}{\omega_k}
\]  

These observations naturally cluster into three distinct regions in the $(A, T)$ plane, where the centroids are the pairs $(A_i, T_i)$ with $i = 1, 2, 3$.
To associate each bin with a particular source, we can apply a \textbf{clustering algorithm} such as \textit{K-means}.  
Once the association is established, we compute the \textbf{binary masks} according to:

\[
M_i(n,\omega_k) =
\begin{cases}
1, & [A(n,\omega_k), T(n,\omega_k)] \ \text{belongs to cluster} \ [A_i, T_i] \\[5pt]
0, & \text{otherwise}
\end{cases}
\]

\begin{figure}[H]
    \centering
    \includegraphics[width=0.4\linewidth]{Data clustering.png}
    \includegraphics[width=0.42\linewidth]{Data clustering 2.png}
    \caption{Examples of data clustering: phase difference $P$ vs level ratio $A$}
\end{figure}

\textbf{Space-time features} should be designed to work synergistically with the clustering algorithm.
For example, \textit{K-means} can be sensitive when one feature has a much larger variance with respect to the others.  
Clustering may also occur in a multi-dimensional space, meaning that each time-frequency bin can be represented with more than two features.  
A suitable choice for time-frequency bin clustering via \textit{K-Means} is to use a feature vector containing \textbf{normalized amplitudes} $A_i$ and \textbf{normalized phase difference} $P$:

\[
A_i(n,\omega_k) = \frac{\abs{Y_i(n,\omega_k)}}{\sqrt{\sum_{m=1}^{M} \abs{Y_m(n,\omega_k)}^2}},
\qquad
P(n,\omega_k) = \frac{1}{2\pi} \angle \left( \frac{Y_2(n,\omega_k)}{Y_1(n,\omega_k)} \right)
\]

The feature vector is defined as follows:

\[
\Phi(n,\omega_k) = 
\begin{bmatrix}
A_1(n,\omega_k) \\
A_2(n,\omega_k) \\
P(n,\omega_k)
\end{bmatrix} \in \mathbb{R}^3
\]

\begin{figure}[H]
    \centering
    \begin{minipage}[c]{0.3\textwidth}
        \centering
        \includegraphics[width=\linewidth]{Mixing masks.png}
    \end{minipage}
    % \hspace{1cm}
    \begin{minipage}[c]{0.3\textwidth}
        \centering
        \includegraphics[width=\linewidth]{Source separation.png}
    \end{minipage}
    % \hspace{1cm}
    \begin{minipage}[c]{0.3\textwidth}
        \centering
        \includegraphics[width=\linewidth]{Source spectogram.png}
    \end{minipage}
    \caption{Example of mixing masks $M_i$ (left), resulting separated spectograms (center) and original sources spectograms $S_i$ (right)}
\end{figure}

% \begin{figure}[H]
%     \centering
%     \includegraphics[width=0.33\linewidth]{Mixing masks.png}
%     \includegraphics[width=0.3\linewidth]{Source separation.png}
%     \includegraphics[width=0.33\linewidth]{Source spectogram.png}
%     \caption{Example of mixing masks $M_i$ (left), resulting separated spectograms (center) and original sources spectograms $S_i$ (right)}
% \end{figure}

\subsection{Source separation: Independent Component Analysis (ICA)}
\textbf{Independent Component Analysis (ICA)} is a technique studied when the number of microphones $M \geq N$ is greater than or equal to the number of sources $N$ (in the following, $M = N = 2$).
It applies in situations where no information is available about the mixing matrix (Blind Source Separation) and where the environment is non-reverberant (instantaneous mixing).

\clearpage

The instantaneous mixing process, in vector form, is expressed as:

\[
\underline{\mathrm{y}} = \left[\mathrm{A}\right] \underline{\mathrm{s}}
\quad \Rightarrow \quad
\underline{\mathrm{s}} = \left[\mathrm{W}\right] \underline{\mathrm{y}}
\]

The goal is to obtain a demixing matrix $\left[\mathrm{W}\right] = \left[\mathrm{A}\right]^{-1}$ without any prior knowledge of the source signals.
The \textbf{fundamental hypothesis} of ICA is that the \textbf{sources} are statistically \textbf{independent}, meaning that their correlation matrix $\left[\mathrm{R}_{\mathrm{s}}\right]$ is diagonal:

\[
\left[\mathrm{R}_{\mathrm{s}}\right] = E\left\{ \underline{\mathrm{s}} \, \underline{\mathrm{s}}^{T} \right\}
\]

Additionally, the source signals have non-Gaussian statistics.

\subsubsection{Ambiguities of ICA}

There are two main \textbf{ambiguities}.
First, we cannot determine \textbf{variances of independent components}: $\underline{\mathrm{s}}$ and $\left[\mathrm{A}\right]$ are unknown, any scalar multiplier in one of the sources can be cancelled by dividing the corresponding column of $\left[\mathrm{W}\right]$ by the same quantity.
Second, We cannot determine the \textbf{order of independent components}: $\underline{\mathrm{s}}$ and $\left[\mathrm{A}\right]$ are unknown, we can freely change the order of the terms in the sum $ \underline{\mathrm{y}} = \left[\mathrm{A}\right] \underline{\mathrm{s}} $
without affecting the final result $\underline{\mathrm{y}}$.

\subsubsection{Statistical concepts and visual representation}

We are going to use the following statistical concepts.
\textbf{Independency}:
\[
p(y_{1}, y_{2}) = p_{1}(y_{1}) \, p_{2}(y_{2})
\]

\textbf{Uncorrelatedness} (covariance is zero):
\[
E\{y_{1}y_{2}\} - E\{y_{1}\} \, E\{y_{2}\} = 0
\]

Independency includes uncorrelatedness.
We also remind the definition of \textbf{orthogonality}:

\[
E\{y_{1}y_{2}\} = 0
\]

Let us analyse an illustration of ICA:

\begin{figure}[H]
    \centering
    \includegraphics[width=0.7\linewidth]{ICA.png}
    \caption{Illustration of ICA}
\end{figure}

\clearpage

On the left, we observe the joint distribution of two independent source signals $\underline{\mathrm{s}}_1$ and $\underline{\mathrm{s}}_2$, uniformly distributed within a square, according to the following probability density function (p.d.f.):

\[
p(s_i) =
\begin{cases}
\frac{1}{2\sqrt{3}}, & \abs{s_i} \leq \sqrt{3} \\[5pt]
0, & \text{otherwise}
\end{cases}
\]

In the center, the mixing process is represented as:

\[
\underline{\mathrm{y}} = \left[\mathrm{A}\right] \underline{\mathrm{s}},
\qquad
\left[\mathrm{A}\right] =
\begin{bmatrix}
2 & 3 \\
2 & 1
\end{bmatrix}
\]

This operation transforms the original independent sources into observed mixtures.

On the right, we see the joint distribution of the mixtures $y_1$ and $y_2$. 
The resulting distribution is an oriented parallelogram, showing that the mixtures are no longer independent.
By identifying the edges of this parallelogram in the $(y_1, y_2)$ space, it is possible to estimate the mixing matrix $\left[\mathrm{A}\right]$, which is a crucial step in ICA for recovering the original independent sources.

\subsubsection{Non-Gaussianity}

The fundamental hypothesis of ICA stated that the \textbf{sources} are statistically independent and have \textbf{non-Gaussian statistics}.
If the independent components were Gaussian, the joint distribution of the mixtures $\underline{\mathrm{y}}$ (in the case of an orthogonal mixing matrix) would be symmetric:

\begin{figure}[H]
    \centering
    \includegraphics[width=0.35\linewidth]{Gaussian joint density.png}
    \caption{Joint density for Gaussian independent components}
\end{figure}

This symmetry means there is no information about the direction of the columns of the mixing matrix $\underline{\mathrm{A}}$.
If the source variables are Gaussian, any orthogonal mixing matrix produces mixed components that are still independent and Gaussian.
In this case, \textbf{source separation is not possible}.

Let us remind of the \textbf{Central Limit Theorem}: the distribution of a sum of independent random variables tends toward a Gaussian distribution.
This property can be exploited to estimate the mixing matrix $\left[\mathrm{A}\right]$. 
Let us consider a generic mixture of the microphone signals:

\[
x = \underline{\mathrm{w}}^T \underline{\mathrm{y}} = \underline{\mathrm{w}}^T \left[\mathrm{A}\right] \underline{\mathrm{s}}
\]

Since $x$ is a linear combination of the components of $\underline{\mathrm{s}}$, it is always more Gaussian than the individual sources themselves.
Therefore, if we choose $\underline{\mathrm{w}}$ such that the \textbf{non-Gaussianity} of $x$ is \textbf{maximized}, we move closer to recovering one of the original sources.  

The vector $\underline{\mathrm{w}}$ does not provide information about the sign of $\underline{\mathrm{s}}$, leading to a \textbf{sign ambiguity}.
Moreover, perfect separation is achieved when $\underline{\mathrm{w}}$ corresponds to one of the rows of the demixing matrix $\left[\mathrm{W}\right] = \left[\mathrm{A}\right]^{-1}$, but the order of the recovered components remains undetermined, which is known as \textbf{order ambiguity}.

\subsubsection{Measures of non-Gaussianity}
Let us define various measurements of non-Gaussianity.
\textbf{Kurtosis} is a statistical measure that is zero for Gaussian variables and non-zero for non-Gaussian ones.  
It is defined as:

\[
\mathrm{kurt}(x) = E\{x^4\} - 3 \left( E\{x^2\} \right)^2
\]

When $\mathrm{kurt}(x) < 0$, the random variable $x$ is classified as \textbf{sub-Gaussian}.
Oppositely, if $\mathrm{kurt}(x) > 0$, the variable $x$ is \textbf{super-Gaussian}.  
A practical approach for quantifying non-Gaussianity is to consider the absolute value $\abs{\mathrm{kurt}(x)}$ as a non-Gaussianity index.

Another measurement involves \textbf{negentropy}.
A key result in information theory states that a Gaussian random variable has the largest entropy among all random variables with the same variance.
The \textbf{entropy} of a random variable $x$ with probability density function $f(x)$ is given by:

\[
H(x) = - \int f(x) \log f(x) \ dx
\]

If we consider a Gaussian random variable $x_{\text{Gaussian}}$ having the same variance as $x$, the negentropy is defined  and can be approximated as follows:  

\[
J(x) = H(x_{\text{Gaussian}}) - H(x) \approx \frac{1}{12} E\{x^3\}^2 + \frac{1}{48} \mathrm{kurt}(x)^2
\]

Maximizing the negentropy is equivalent to maximizing the non-Gaussianity of the signal.

We can consider other measures of independency.
The \textbf{mutual information} quantifies the dependence between random variables.  
For $N$ random variables $x_1, x_2, \dots, x_N$, it is defined as:

\[
I(x_1, x_2, \dots, x_N) = \sum_{i=1}^{N} H(x_i) - H(\underline{\mathrm{x}}),
\qquad H(\cdot): \text{entropy}
\]

Mutual information is always non-negative and equals zero if and only if the variables are statistically independent.  

The \textbf{Kullback–Leibler divergence} is one of the most natural measures of independency: it is defined as the distance between the joint density $f(\underline{\mathrm{x}})$ and the product of its marginal densities.
In other words, the Kullback–Leibler divergence measures how much one probability distribution differs from a second, reference probability distribution.

\subsubsection{Estimation of the demixing matrix}
Once we choose a suitable cost function $J(\underline{\mathrm{w}})$ (such as kurtosis, negentropy, or mutual information), the estimation of the demixing model is formulated as one of the following \textbf{optimization problems}, depending on the selected independence metric: 

\[
\hat{\underline{\mathrm{w}}} = \underset{\underline{\mathrm{w}}}{\arg \max} \ J(\underline{\mathrm{w}})
\qquad \text{or} \qquad
\hat{\underline{\mathrm{w}}} = \underset{\underline{\mathrm{w}}}{\arg \max} \ J(\underline{\mathrm{w}})
\]

\clearpage

In practice, the demixing model is usually estimated through LMS-like \textbf{iterative algorithms}.
Since the \textbf{cost function} is typically highly \textbf{non-linear}, these algorithms can be trapped in \textbf{local minima or maxima}.
As a consequence, the obtained solution can depend on the \textbf{initial conditions} used for each realization:

\begin{figure}[H]
    \centering
    \includegraphics[width=0.3\linewidth]{ICA convergence.png}
    \caption{Demixing matrix estimation: example of convergence starting from different initial conditions}
\end{figure}

\clearpage
